% !PRINTED 138
\weberpage{138}
%<<CONT
are proper divisors of these numbers, have already been proved to be metacyclic.

There remains only the investigation of the numbers 60 and 90. That a simple group of degree 60 exists we know already; namely the alternating group of the permutations of five letters (Bd.\,I, \S.\,177). But it is still questionable whether this is the only one.

A simple group $P$ of degree 60 contains, by \S.\,29, I., as divisor a group of degree 5 and a group of degree 3. If in the second \gesperrt{Sylow's} theorem (\S.\,30) one takes for $Q$ the group of degree 5, then $R$ must be of degree 10, because the index of $R$, being congruent to 1 with respect to the modulus 5, can only be $= 6$. Hence by \S.\,28, 2. the group $P$ of degree 60 can be represented as a transitive permutation group of six figures; but it cannot contain any cyclic permutation of only three figures, because as a simple group it must (by Bd.\,I, \S.\,158, 2.) be primitive, and therefore, if it contained a three-term cycle, it would by Bd.\,I, \S.\,153, 10. have to contain the whole alternating group of degree 360. We may therefore assume one of the permutations of degree 3 in the form:
\[
a = (1, 2, 3) (4, 5, 6).
\]
To this we will add a permutation of degree 5, $b$, which can only be a cyclic one. If we let the figure 6 be absent from it, then in place of $b$ we can take such a power of $b$ (which must of course also occur in $P$) that, say, 1, 2 become the first two figures, and then six possibilities remain for $b$:
\[
(1, 2, 3, 4, 5), \quad (1, 2, 4, 3, 5), \quad (1, 2, 4, 5, 3),
\]
\[
(1, 2, 3, 5, 4), \quad (1, 2, 5, 3, 4), \quad (1, 2, 5, 4, 3);
\]
of these six assumptions, however, the three standing in one row are again not essentially different (i.\,e. different only in the designation). For if, for example, in the second assumption $b = (1, 2, 4, 3, 5)$ one replaces the element $a$ by $a^2$ and $b$ by $b^3$, one obtains
\[
(1, 3, 2) (4, 6, 5); \quad (1, 3, 2, 5, 4),
\]
which, by the permutation of the figures 2 with 3 and 4 with 5, passes into the first assumption. The third assumption $b = (1, 2, 4, 5, 3)$ passes, if one replaces $b$ by $b^4$ and $a$ by $a^2$ and then
%CONT>>
